% !TeX root = main.tex
\lecture{8}{Thu 26 Feb 2026 14:00}{Cosmology II: Expansion of the Universe and Key Principles}

\section{Expansion of the Universe}
Data from supernovae found that the universe was not expanding at a constant rate. The rate of expansion is increasing and accelerating.

The best current model we have is the Hot Big Bang Model:
\begin{figure}[H]
    \centering
    \includegraphics[width=0.75\textwidth]{figures/lec08-01.png}
     \caption{}
\end{figure}

This is where the universe was created from an initial singularity, resulting in the big bang - an explosion.

Space itself is expanding and the universe is shaped by gravity as a fundamental force. The fundamental particles were created first in a dense and hot universe. As the universe expanded, it cooled down allowing for the propagation of CMB photons and the creation of meaningful cosmic structure.

\subsection{Big Bang Cosmology}
For the rest of the course we will:
\begin{itemize}
    \item Explore big-bang centric cosmology, in the context of Newtonian gravity.
    \item We'll build and discuss the main equations that model the evolution of the universe. Using these we'll explore:
    \begin{itemize}
        \item The history of the universe.
        \item The composition of the universe.
        \item The eventual fate of the universe.
    \end{itemize}
    \item We'll discuss key observations and any problems where these do not fit the models we're applying them to.
\end{itemize}

\section{Key Principles}
\subsection{The Cosmological Principle}
The distribution of matter is:
\begin{itemize}
    \item \textbf{Homogeneous}, meaning that the universe looks the same at every point. There is no special point in the universe, and the average matter density is the same everywhere.
    \item \textbf{Isotropic}, meaning the universe looks the same in all directions
\end{itemize}

However this isn't perfectly fulfilled in our universe. It's a valid approximation on sufficiently large scales $>100$Mpc, but it isn't perfect.

\subsection{Distance Measures - The Parsec}

The parsec is a unit of distance useful on cosmological scales:
\[
    1 \si{pc} = 3.086 \times 10^{16} \si{m}
\]

Suppose an observer on Earth measures a very distant star. We take two measurements six months apart. As the Earth is moving relative to the star, it appears to us that the star has shifted position (in the Earth's rest frame). The total observed angular shift is $2p$, where $p$ is the parallax angle --- the angle at the star in the right triangle formed by the Sun, Earth, and star, with 1\,AU as the opposite side.

The parsec is defined as the distance at which $p = 1$ arcsecond ($4.848 \times 10^{-6}$\,rad), giving the relationship:

\[
    d\,[\si{pc}] = \frac{1}{p\,[\text{arcsec}]}
\]

\subsection{Redshift}
This has already been covered in other modules, but to recap the fractional shift is:
\[
    z = \frac{\lambda_0 - \lambda_e}{\lambda_e} \approx \frac{v}{c}
\]
for non-relativistic sources where $v \ll c$.

\section{Expansion of the Universe}
Hubble found that there was a linear relationship between the recessional velocity of a galaxy and its distance.

\[
    \pmb{v} = H_0 \pmb{r}
\]

Where $H \approx 68 \si{\kilo\metre \s^{-1} Mpc^{-1}}$. Note that the parameter $H$ varies with time. $H_0$ represents the value at the current epoch.

This isn't exact, and does not hold for nearby galaxies where other forces dominate over the expansion of the universe. These galaxies have some peculiar velocity which forms the total velocity:
\[
    v_\text{total} = v_\text{pec} + v_H
\]

For distant galaxies, the recessional velocity dominates. This is called Hubble flow.

\subsection{Comoving Coordinates}
As the universe expands, the physical (proper) distance between two galaxies grows. This means that, in a coordinate system centred on the earth, the coordinates of other galaxies are increasing.

Instead, we define a new coordinate system where the coordinate system itself expands alongside the universe. This means that galaxies remain at fixed coordinates as the universe expands, and the coordinate grid expands at the same rate as te universe. The distance between galaxies in these coordinates is therefore constant.

This rate of increase is given by the universal expansion rate (a scale factor) $a(t)$. Due to homogeneity, this scale factor is only time-dependant, as the rate of expansion is the same everywhere in the universe.

This means that we can relate proper distance $\pmb{r}(t)$ and the fixed comoving distance $x$ with:
\[
    \pmb{r}(t) = a(t) \pmb{x}
\]

